\cotefile{157-2}
\scanpage{259}
Catégorie \unsure{des|du} chemins (2)

$\overline{\text{\uline{VII}}}$

% archive sheet number: 258

\scanpage{260}
\mspage{1}

\unsure{Catégories|Catégorie} de Chemins (nouvelle version)

Comme le formalisme ``simplifié'' des chemins que j'ai voulu dégager à la va-vite s'est révélé d'une lourdeur insupportable, de nature à le rendre quasiment inutilisable, je vais essayer de dégager un formalisme souple, collant à l'intuition, au prix seulement d'un effort conceptuel supplémentaire.

1. \uline{Intervalles} \uline{et \unsure{leur|leurs} \unsure{orientations|orientation}}.

On appelle \uline{catégorie} \uline{intervalle} \uline{\unsure{standard|standart}} une catégorie définie par \deleted{\unsure{un}} \unsure{ens.} ordonné, dont l'ens. des \deleted{\unsure{s}}éléments est un intervalle d'entiers $[0,n]$ ($n\in\mathbb{N}$, on n'exclut pas le cas $n=0$, i.e. de la catégorie ponctuelle $\{0\}$) :

\[\operatorname{Ob}I=[0,n]\subset\mathbb{N},\]

\deleted{\unsure{b)}} et satisfaisant l'hypothèse suivante : a) pour toute \deleted{\unsure{$i\in N$}, \unsure{l'un}} \inserted{paire de} \deleted{des sommets de $I$} deux sommets consécutifs $i-1,i\deleted{\unsure{+1}}$ de $I$ ($1\leq i\leq n\deleted{\unsure{+1}}$), il y a une flèche qui les relie, donc soit $i\longrightarrow i+1$, soit $i\longleftarrow i\unsure{+1|\text{et}}\deleted{\unsure{,\ \ldots}}$) Cette flèche est notée $\underline{u}_i$

\diagram{\begin{tikzcd}[ampersand replacement=\&,column sep=large]
(i-1)\arrow[r,"{\underline{u}_i}"] \& i \& \text{ou} \& (i-1) \& i\arrow[l,"{\underline{u}_i}"'] \& {;}
\end{tikzcd}}

b) La catégorie $I$ est engendrée par l'ens. de ses flèches $\underline{u}_i$ (appelées \uline{flèches de \unsure{transition}}.)

\deleted{Donc il y a au} La flèche \deleted{\unsure{élé}} \unsure{de transition} est dite \uline{directe} si elle va de $i-1$ à $i$, \uline{rétrograde} dans le sens inverse. Si on affecte l'intervalle $[i-1,i]$ du signe $+$ ou $-$ suivant que $\underline{u}_i$ est directe ou rétrograde, on voit donc que les

% archive sheet number: 259

\scanpage{261}
\mspage{2}

intervalles standard de longueur donnée $n$ correspondent biunivoquement aux syst. de signes \deleted{\unsure{$\pm$}} $+$ \inserted{ou $-$} sur les $n$ intervalles de transition $[i-1,i]$ ($1\leq i\leq n$), il y en a donc $2^n$. Le sommet $0$ est l'\uline{origine} \inserted{ou la \uline{source}} de l'intervalle, le sommet $n$ l'\uline{extrémité} ou le \uline{but} de l'intervalle.

Un intervalle standard est \uline{orienté}, du fait qu'on a la relation d'ordre \deleted{\uline{totale}} naturelle sur \deleted{\unsure{l}}l'ens. de ses sommets, \inserted{savoir celle des entiers de $[0,n]$.} Quand on parlera de la relation d'ordre sur $I$, qu'un sommet $i$ en ``\uline{précède}'' ou ``\uline{suit}'' un autre $j$ ($\unsure{i|v}<j$ ou $i\unsure{>|{\geq}}j$ respectivement), c'est de cette relation d'ordre qu'on voudra parler. Celle associée à la structure de catégorie ne sera \unsure{jamais} regardée comme une relation d'ordre.

Il y a également une relation d'ordre totale naturelle sur l'ens. des flèches de transition $\underline{u}_i$ \inserted{$(1\leq i\leq n)$}, donnée par \deleted{\unsure{la}} la relation d'ordre des entiers $i$, $1\leq i\leq n$, qui les indexent. On dira qu'une flèche $\unsure{i|\prime}$ de tr. \uline{précède} ou \uline{suit} une autre, ou \deleted{\unsure{la}} qu'elle est \uline{successeur} \inserted{\uline{immédiat}} ou \uline{prédécesseur immédiat} d'une autre (si $i=j+1$, ou $i=j-1$).

Si on a un ensemble d'entiers

\[E\subset\mathbb{Z},\]

on \deleted{\unsure{dira}} considérera dans $E$ la relation d'équivalence engendrée par la relation entre deux éléments $(i,j)$ de $E$ : $j$ est successeur immédiat de $i$ \uline{dans} $\deleted{\unsure{j}}\mathbb{Z}$ (i.e. $j=i+1$). Deux éléments \inserted{distincts $i,j$} de $E$ \deleted{qui} sont congrus pour cette relation \unsure{sss|ssi} \deleted{\unsure{on a}} l'intervalle d'entiers qu'ils déterminent est $\subset E$, donc

\[\text{si }i<j,\quad i\text{ et }j\text{ sont ``\uline{connectés}'' dans }E\text{ ssi }[i,j]\subset E.\]

% archive sheet number: 260

\scanpage{262}
\mspage{3}

Quand on parlera des ``composantes connexes'' de $E$, il s'agira toujours des classes suivant cette relation d'équivalence.

Nous \unsure{identifions|identifierons} l'ens des flèches de transition \inserted{$\mathrm{Fl.tr.}(I)$} de $I$ à l'intervalle d'entiers $[1,n]$, donc à une partie de $\mathbb{Z}$. Si on a une partie $E\subset\mathrm{Fl.tr.}(I)$, \deleted{\unsure{elles sont}} \inserted{on y a} donc \deleted{\unsure{aussi}} la notion de flèches de transition ``connectées'' dans $E$, et la notion de connexité pour $E$.

Si $i<j$, $\underline{u}_i$ et $\underline{u}_j$ sont connectées \unsure{sss|ssi} \unsure{toutes} les flèches de transition intermédiaires entre $\underline{u}_i$ et $\underline{u}_j$ sont dans $E$. Et $E$ est connexe sss il est soit vide, soit formé d'une suite de flèches de transition dont chacune est successeur immédiat de celle qui \unsure{la|le} précède

\[E=\{\underline{u}_i,u_{i+1},\ldots u_{i+l-1}\},\]

où $l$ est le \unsure{``}cardinal de $E$. Une telle $E$ est appelée une \uline{\unsure{chaine|chaîne}} de flèches de transition, et $l$ est sa longueur.

Une chaîne est dite \uline{composable} si \unsure{le|la} \deleted{flèche $\unsure{u}_{i+l-1}$ \illegible{1}} \deleted{\unsure{suite}} \inserted{composé $\underline{u}_{i+l-1}\circ\cdots\underline{u}_{i+1}\circ\underline{u}_i$ a un sens, i.e. si elles sont toutes de $\mathord{=}$ sens, soit direct, soit rétrograde.} On désigne par $\mathrm{Fl.tr.}^{+}(I)$, $\mathrm{Fl.tr.}^{-}(I)$ \deleted{les deux} \inserted{l'ens.} des fl. de tr. directes, ou rétrogrades, respectivement.

Leurs \deleted{\unsure{con}} composantes connexes ne sont autres que les chaînes composables maximales. \deleted{\unsure{qu'on appellera}} \inserted{Ainsi, les} \deleted{les \uline{chaînes} \uline{fondamentales}} chaînes composables maximales forment une partition de l'ens. des \deleted{chaînes} \inserted{flèches de} \unsure{et|de} transition, en intervalles d'entiers mutuellement disjoints. Ces chaînes forment elles-$\mathord{=}$s un ens. totalement ordonné

\[C_1,C_2,\ldots,C_r\qquad(1\leq r\leq n)\]

L'entier $r$ s'appelle la ``\uline{longueur utile}'' de l'intervalle

% archive sheet number: 261

\scanpage{263}
\mspage{4}

$I$, c'est donc le cardinal de l'ens \unsure{des|de} ses chaînes composables maximales.

Deux chaînes composables maximales sont toujours de \unsure{signe|signes} \unsure{opposé|opposés}, si l'une est directe l'autre est rétrograde, et inversement. \deleted{Si la pre} La première chaîne composable maximale \inserted{$C_1$} (disons simplement ``\uline{chaîne structurale}'') a le $\mathord{=}$ signe que $u_1$ ; \deleted{s'il est $+$} \deleted{i.e. si $u_1$ est directe,} on dit que la \uline{direction initiale} de $I$ est directe, resp. rétrograde, s'il en est ainsi de $u_1$, ou ce qui revient au $\mathord{=}$, de $C_1$.

De $\mathord{=}$ on définit la \uline{direction finale} de $I$ comme celle de $u_n$, ou encore celle de $C_r$. \deleted{Ce n'est que dans le} Direction initiale et finale sont égales si la longueur utile $r$ de $I$ est \inserted{im}paire, inégales si elle est paire. (\uline{\unsure{NB}} Ce n'est que pour $n=0$, i.e. quand $I=\{0\}$ est la cat. ponctuelle, que direction initiale et direction finale ne sont pas définies).

Les chaînes composables correspondent aux sous-catégories pleines de $I$ qui sont isomorphes à une catégorie $\mathbb{\Delta}_i$ (où $i$ est la longueur de la chaîne). Notons que $\mathbb{\Delta}_i$ n'a pas d'autre automorphisme que l'identité, donc les chaînes maximales \inserted{de longueur donnée $i\geq1$} correspondent exactement aux foncteurs

\diagram{\begin{tikzcd}[ampersand replacement=\&]
\mathbb{\Delta}_i\arrow[r,hook] \& I
\end{tikzcd}}

qui sont \deleted{\unsure{des mono}} injectifs sur les sommets. Un tel foncteur est \deleted{alors \unsure{automatiq}} pleinement fidèle, et un monomorphisme. \deleted{Une chaîne} Ce foncteur définit une application \inserted{\unsure{str|str.}} croissante ou \deleted{\unsure{détr.}} décroissante sur

% archive sheet number: 262

\scanpage{264}
\mspage{5}

les sommets, suivant que la chaîne est directe ou rétrograde.

Notons que tout \inserted{simplexe type $\unsure{i|j}\geq0$} $\mathbb{\Delta}_i$ peut être vu comme un intervalle, de longueur $i$, de longueur utile $1$ (si $i\geq1$, $0$ si $i=0$). Son origine est $0$, son extrémité est $i$.

Mais appelons une catégorie quelconque une \deleted{intervalle} \uline{catégorie} \uline{intervalle} si elle est isomorphe à une \deleted{catégorie} intervalle standard. Notons ici

\uline{Proposition 1} Soit $I$ un intervalle \unsure{standart|standard} de longueur $n\geq1$. \deleted{Alors origine et extrémité \unsure{coïncident|coïncidant}}

a) \deleted{Considérons} Soit $i\in I$. Pour que \deleted{$i\in\{0,n\}$ i.e.} \inserted{$i\in\{0,n\}$, i.e.} pour que $i$ soit un des deux ``\uline{bouts}'' \inserted{(ou ``\uline{extrémités}'')} (source ou but) de $I$, il f. et s. \unsure{ou|que} l'une des deux catégories $I_{/i}$ ou $i\backslash I$ soit \deleted{\unsure{triviale}} la catégorie ponctuelle \inserted{et que l'autre soit iso. à un $\mathbb{\Delta}_j$ $*$} (Caractérisation intrinsèque des bouts d'un intervalle.) Pour que $i$ soit la source, il f. et s. que \unsure{$i$ est}

% slanted
\margin{$*$ où $j$ est la longueur de l'unique chaîne structurale qui ``touche'' $i$ (i.e. telle que $i$ soit \deleted{dans} \unsure{la|$\in$} sous-catégorie pleine \unsure{''} définie par ladite chaîne)}

b) Tout foncteur $I\longrightarrow I$ qui fixe les deux \unsure{extrémités} de $I$ est l'identité.

\deleted{c) \unsure{Il en}}

\uline{Corollaire} Le groupe des automorphismes de $I$ a au plus deux éléments. S'il en a deux, l'autom. non trivial échange les deux extrémités. (On dit \deleted{\unsure{que}} alors que l'intervalle est \uline{symétrique}.) \deleted{\unsure{on a trivial}}

\deleted{\uline{Proposition 2}. Si $I$ est un intervalle standard \deleted{\unsure{$I$}} il existe un unique couple $(I',\Theta)$, \unsure{avec|où} $I'$ un \deleted{est isomorphe à un intervalle standard unique} intervalle standard, et $\Theta$ un isomorphisme de $I$ sur $I'$, qui transforme but en origine et origine en but)

$I'$ est déduit de $I$, en transformant, par la symétrie}

\[\deleted{i\longmapsto n-i\qquad [0,n]\longrightarrow[0,n]\,,}\]

% archive sheet number: 263

\scanpage{265}
\mspage{6}

\deleted{le système de signes \uline{opposé} de celui de $I$. \inserted{\unsure{Je vais expliciter}}}

NB Si $i$ n'est pas un bout, et si néanmoins l'une des deux catégories $i\backslash I$, $I_{/i}$ est ponctuelle, disons $I_{/i}$, alors au voisinage immédiat de $i$ $I$ a la forme

\diagram{\begin{tikzcd}[ampersand replacement=\&,column sep=large]
\underset{i-1}{\vphantom{\bullet}} \& \underset{i}{\bullet}\arrow[l,"{\underline{u}_{\deleted{\unsure{i-1|i+1}}\unsure{i}}}"']\arrow[r,"{\underline{u}_{i+1}}"] \& \underset{i+1}{\vphantom{\bullet}}
\end{tikzcd}}

\deleted{(\unsure{$u_i$ et $u_{i+1}$}} ($u_i$ rétrograde, $u_{i+1}$ directe -- si c'était $i\backslash I$ la ponctuelle, alors $u_i$ serait directe, $u_{i+1}$ rétrograde -- en tous les cas \deleted{\unsure{deux}} \unsure{ces} deux flèches sont de sens opposé, i.e. non composables). \deleted{\unsure{alors}} On trouve donc que $i\backslash I$ est une catégorie de la forme

\diagram{\begin{tikzcd}[ampersand replacement=\&,row sep=small,column sep=small]
 \& \& i\arrow[dl]\arrow[dr] \& \& \\
 \& i-1\arrow[ddl,dotted,no head]\arrow[ddl,shorten <=1.4em] \& \& i+1\arrow[dr,dotted,no head]\arrow[dr,shorten <=1.1em] \& \\
 \& \& \& \& i+\unsure{s} \\
 i-r \& \& \& \&
\end{tikzcd}}

où $r$ est la longueur \deleted{\unsure{entière}} de la chaîne structurale contenant $\underline{u}_i$, $s$ celle de la chaîne structurale contenant $\underline{u}_{i+1}$. Cette catégorie est une somme amalgamée $\mathbb{\Delta}_r\amalg\mathbb{\Delta}_s$ (recollées suivant leur origine), mais n'est pas isom. à un $\mathbb{\Delta}_j$, p.ex. faute d'avoir un objet final.

Cette observation prouve a). Pour prouver b, il suffit de prouver que l'ordre total sur $I$ est déterminé de façon intrinsèque, \deleted{\unsure{On va voir}} \deleted{\unsure{Pour|Donc} ceci} par \unsure{le|les} choix de l'\uline{origine} \inserted{$s_0$} de $I$ (parmi ces deux extrémités). Supposons que l'on ait déjà déterminé intrinsèquement les sommets $s_0,s_1,\ldots s_i$ ($i\unsure{\geq|>}0$). Comment déterminer $s_{i+1}$ ? \deleted{(On regarde \unsure{$s_i$})} \inserted{(quand $i<n$)} Soit \deleted{$J=\{s_0,\ldots s_{i-1}\}$} $I'$ la sous-catégorie pleine de $I$

% archive sheet number: 264

\scanpage{266}
\mspage{7}

formée des $j\in I$ qui ne sont pas parmi les $s_j$ déjà construits, sauf \deleted{\unsure{du}} le dernier $s_i$. \deleted{\unsure{Cet $i$}, \unsure{$I'$}} \inserted{(donc $s_i\in I'$, $\operatorname{card}I'\geq2$)} Alors \deleted{l'une des catégories} $I'$ est isom. à un intervalle \unsure{standard|standart}, \unsure{son|avec} $s_i$ correspondant à l'origine. Une des catégories $i\backslash I$, $I_{/i}$ est ponctuelle, l'autre \deleted{\unsure{équival}} isom. à une catégorie $\mathbb{\Delta}_j$, et de plus,

$i$ est objet initial ou final de $\mathbb{\Delta}_j$. Si c'est l'objet initial $0$, on prend pour $s_{i+1}$ le successeur \inserted{immédiat} de $i$ dans $\mathbb{\Delta}_j$. Si c'est l'objet final $j$, on prend pour $s_{i+1}$ le prédécesseur immédiat $\unsure{j'|j}-1$.

Si on était parti du \unsure{but|bout} \inserted{$\unsure{b|h}$} de $I$, au lieu de partir de la source, on aurait trouvé la relation d'ordre totale opposée. \deleted{Pour cela} \unsure{Affectant|Affectons} l'ancien but du numéro d'ordre $0$, on trouve une renumérotation de $I$, $i$ devenant $n-i$, et les \deleted{\unsure{deux}} flèches de transition qui étaient \deleted{d'un signe} directes pour la \deleted{première} relation d'ordre initiale, deviennent rétrogrades dans la relation \deleted{\unsure{oppos}} \inserted{d'ordre} opposée, et inversement. On voit ainsi :

\uline{Proposition 2} Si $I$ est un intervalle \inserted{standard}, il existe un unique \deleted{intervalle $I'$,} couple $(I',\Theta)$, où $I'$ est un intervalle standard, et $\Theta$ un iso $I\simeq I'$ qui transforme le but de $I$ dans la source de $I'$ (et la source de $I$ en le but de $I'$.

Cet intervalle $I'$ est appelé l'intervalle \uline{opposé} de $I$. Il faut le voir, intrinsèquement, comme la ``même'' catégorie $I$, mais avec une ``orientation'' opposée.

% archive sheet number: 265

\scanpage{267}
\mspage{8}

Une façon plus claire de voir ceci est la suivante. Soit $J\deleted{\unsure{*}}$ une catégorie \inserted{intervalle, i.e.} isomorphe à un intervalle \inserted{standard}. Alors il y a exactement \inserted{(si $I\ne e$)} \uline{deux} couples $(I,\Theta)$ d'un intervalle standard $I$ et d'un isom. $\Theta:I\xrightarrow{\sim}J$.

Ces deux isomorphismes induisent sur $J$ deux ordres \deleted{\unsure{tot}} totaux \deleted{inverses} \inserted{opposés} l'un de l'autre, \deleted{et sont} Chaque couple est déterminé de façon unique par ledit ordre \deleted{total}, qu'on appellera une \uline{orientation} de l'intervalle $J$.

% slanted
\margin{\uline{NB} On parlera de l'intervalle \uline{orienté} \uline{opposé} $I^{\mathrm{opp}}$ d'un intervalle \uline{orienté} $I$, qu'il faut se garder de confondre avec la catégorie opposée $I^{\unsure{\circ|\mathrm{o}}}$. $I^{\mathrm{opp}}$ a $\mathord{=}$ catégorie \unsure{ss-jacente} que $I$, mais avec une orientation opposée.}

Ainsi $J$ a deux orientations, et \deleted{pour} chacune détermine un \uline{isomorphisme} \uline{orienté} (i.e. respectant les orientations) de $J$ sur un intervalle standard. \unsure{Et|Si} les deux intervalles sont les mêmes, ssi. il existe un automorphisme de $J$ \inserted{(nécessairement unique)} qui échange ses deux orientations, ou ce qui revient au $\mathord{=}$, qui échange entre elles les deux extrémités. (Ainsi les catégories opposées l'une de l'autre

\diagram{\begin{tikzcd}[ampersand replacement=\&,column sep=small,row sep=small]
 \& \cdot\arrow[dl]\arrow[dr] \& \& \& \cdot \& \\
 \cdot \& \& \cdot \& \cdot\arrow[ur] \& \& \cdot\arrow[ul] \& {,}
\end{tikzcd}}

sont chacune des intervalles symétriques.) \inserted{Choisir une orientation revient \unsure{ainsi|aussi} à choisir l'origine parmi les deux extrémités.}

Notons que la \uline{catégorie opposée} à un intervalle est un intervalle, \deleted{que} et la cat. opposée à un intervalle standard \deleted{\unsure{n'}} est \deleted{\unsure{auto}} l'intervalle standard qui s'en déduit en changeant tous les signes sur les intervalles de transition $[i-1,i\deleted{\unsure{-1}}]$ ($1\leq i\leq n$). On se gardera de confondre la propriété $I\simeq I^{\circ}$, avec la propriété de symétrie.

% archive sheet number: 266

\scanpage{268}
\mspage{9}

Ainsi les deux catégories ($\unsure{\Psi|\Phi}_1$ et $\Phi_1$) ci-dessus sont symétriques, mais opposées l'une de l'autre sans être isomorphes. D'autre part \deleted{$\mathbb{\Delta}$} les $\mathbb{\Delta}_i$ ne sont pas symétriques, mais \unsure{chacuns|chacun|chacune} \unsure{isomorphes|isomorphe} à la catégorie opposée.

Si $I$ est isomorphe à la catégorie opposée $I^{\unsure{\circ|\mathrm{o}}}$, et si $\Theta$ est un tel isomorphisme (\unsure{néc.} \deleted{\unsure{pr}}unique), alors \deleted{\unsure{il sera}} \deleted{il est immédiat} il est immédiat que $\Theta$ renverse l'orientation (sinon ce serait l'identité sur les sommets, et les flèches de transition seraient à la fois directes et indirectes), \deleted{de} donc $\Theta(0)=n$, $\Theta(1)=n-1$, et on voit que $u_1$ et $u_n$ sont de même signe. Mais si $I$ est symétrique, $u_1$ et $u_n$ sont de signes opposés. Donc $I$ ne peut être à la fois symétrique, et isomorphe à \deleted{\unsure{$I^\circ$}} l'intervalle opposé.

Pour toute catégorie $\mathbb{\Delta}_i$, celle-ci étant un intervalle, il y a \inserted{(si $i>0$)} \uline{deux orientations}. \deleted{L'orientation} En tant qu'intervalles orientés, il y a $\mathbb{\Delta}_i^+$ (avec l'\deleted{relation}orientation ordinaire, donnée par l'ordre ordinaire sur $\operatorname{Ob}\mathbb{\Delta}_i=[0,i]$), et $\mathbb{\Delta}_i^-$ avec l'orientation opposée.

\deleted{\unsure{Ce|On}}

\deleted{\uline{Proposition} Soit $I$ une \deleted{\unsure{inter}} catégorie intervalle orientée, (de longueur $n\geq1$) donc can. isom. à un intervalle standart. Alors il y a une suite unique de}

\[\deleted{\mathbb{\Delta}_{l_1},\ldots,\mathbb{\Delta}_{l_r}}\]

\deleted{de simplexes types, \deleted{de long} de dimensions $l_1,\ldots l_r>0$, \deleted{chacun orienté} et une suite unique d'orientations de ceux-ci, de telle façon que}

% archive sheet number: 267

\scanpage{269}
\mspage{10}

2. \deleted{\uline{\unsure{Composition} des intervalles orientés}} \inserted{\uline{Sous-catégories pleines d'un intervalle}}

Soit $I$ une \inserted{catégorie} intervalle orientée, \inserted{de longueur $n\geq1$. Nous identifions $\operatorname{Ob}I$ à l'ens d'entiers $[0,n]$.} Pour une chaîne \deleted{de} \inserted{$\underline{u}_i,\underline{u}_{i+1},\ldots,\underline{u}_{i+l-1}$} \unsure{flèches} de transition de $I$, on appelle \uline{domaine} de la chaîne, l'ens des sources et des buts de ces éléments. C'est donc un intervalle d'entiers, $[i-1,i+l-1]$, de longueur \inserted{$l$} égale à celle de la chaîne, dont l'origine est déterminée par la première flèche $u_i$ de la chaîne (c'est sa source, si $u_i$ est directe, et son but si $u_i$ est rétrograde), et dont l'extrémité est déterminée par la dernière $u_{i+l-1}$.

Les intervalles d'entiers \inserted{(vus comme des sommets de $I$)} \unsure{correspondant|correspondants} aux chaînes structurales sont appelés les \uline{intervalles structuraux} \inserted{ou \uline{intervalles critiques}} de $I$. Contrairement aux chaînes structurales, ces intervalles ne forment pas une partition de $I$, puisque deux intervalles structuraux peuvent avoir une extrémité commune. Si les chaînes structurales sont

\[C_1,C_2\ldots,C_r\qquad (r=\text{long. utile de }I)\]

de longueurs

\[l_1,l_2,\ldots,l_r,\]

alors les intervalles structuraux sont

\[\underbrace{[0,l_1]}_{I_1}\underbrace{[l_1,l_1+l_2]}_{I_2}\unsure{,}\ldots\ \underbrace{[l_1+\cdots+l_{r-1},l_1+\cdots+l_{r-1}+l_r]}_{I_r},\]

une suite d'intervalles d'entiers dont deux consécutifs ont une extrémité commune, et tels que, si $i<j$, chaque élément de $\unsure{I|I_i}$ soit $\leq$ chaque élément de $I_j$. On appelle \uline{sommets critiques} de $I$ les extrémités des intervalles structuraux.

\uline{Prop. \deleted{\unsure{2}}3} Soit $s\in I$, \inserted{soit $I'=\deleted{\unsure{{}_{s}\backslash}}I_{/s}$ (sous-cat. pl. de $I$ des $s'$ tels que existe $s'\longrightarrow s$), $I''=I_{/s}$, $I'_0=I_{/s}\backslash\{s\}$, $I''_0=\unsure{{}_{s}\backslash I|I_{/s}}\backslash\{s\}$.} \deleted{Pour que $s$ soit critique, il f. et s. que les deux sous-catégories} \deleted{Pour que} Pour que $s$ soit critique, il f. et s. que $I\backslash\{s\}$ soit cat. somme des $I'_0$ et $I''_0$ ; \deleted{et que celle au moins \unsure{un vide}}

% slanted
\margin{i.e. qu'on ait : \unsure{a)} Il n'y a pas de flèche entre un élément de $I'_0$ et un élément de $I''_0$. \deleted{b) $I'_0$ et $I''_0$ sont \unsure{pleines} \illegible{1}} (\unsure{évidemment} \illegible{1} par \unsure{la nécessité}.)}

% archive sheet number: 268

\scanpage{270}
\mspage{11}

\deleted{Soient} \uline{Proposition 4} Soient \deleted{$i<j$} \inserted{$i,j$} dans $I$, \inserted{avec $i<j$.} Conditions équivalentes

a) Il existe une flèche entre $i$ et $j$ (soit $i\to j$, soit $j\to i$)

b) $i$ et $j$ appartiennent à un même intervalle structural \inserted{(\unsure{néc\textsuperscript{t}} unique).}

c) La sous-catégorie pleine de $I$ définie par l'ens d'objets $[i,j]$ est isomorphe à un $\mathbb{\Delta}_{\alpha}$ (on aura $\alpha=j-i$)

On dira alors que $i$ et $j$ sont \uline{directement connectés} dans $I$. (Et on utilisera la $\mathord{=}$ terminologie dans une cat. quelconque : $i$ et $j$ sont dir. conn. \unsure{sss|ssi} $\operatorname{Hom}_I(i,j)$ \uline{ou} $\operatorname{Hom}_I(j,i)$ est non vide.)

\uline{Proposition} 5. Soit $J\subset I$, \deleted{\unsure{consi}} une partie de $I$, considérons la sous-catégorie pleine de $I$ qu'elle définit. Alors

a) Les composantes connexes \inserted{$J_\alpha$} de $J$ sont des intervalles, \deleted{(orientés par l'ordre} et l'ordre induit sur \inserted{un} $J_\alpha$ par celui de $I$ est une orientation (\deleted{\unsure{dite}} \uline{orientation induite}) de $J_\alpha$.

% slanted
\margin{On appelle \uline{sous-intervalle} de $I$, une sous-catégorie pleine connexe, i.e. qui est \unsure{elle-même|elle-$\mathord{=}$} un intervalle. Notons d'ailleurs qu'une sous-catégorie $I'$ de $I$ qui est un intervalle, est nécessairement pleine, cf plus bas \S{} des foncteurs d'intervalles.}

b) \inserted{Supposons \deleted{\unsure{$J\ne\varnothing$}} $\operatorname{card}J\geq1$.} Soit $i$ le plus petit élément de $J$, $j$ le plus grand élément, \deleted{Supposons} donc $i\leq j$. Pour que $J$ lui-même soit un intervalle, il f. et s. que tout sommet critique contenu dans $[i,j]$ soit élément de $J$.

\uline{NB} Si $s$ \inserted{$\in[i,j]$} est un sommet critique \deleted{de $[i,j]$ \unsure{et}} \inserted{de $I$}, \deleted{Si $i,j$ sont deux éléments de $I$ qui} \deleted{qui n'est pas dans $J$, alors} qui n'est pas \deleted{\unsure{dans|en}} $J$, alors par prop. 3 il disconnecterait $J$ \deleted{\unsure{en} $\unsure{2|q}$}, i.e. $J$ serait somme directe de deux ss-catégories $J'$ et $J''$ non vides ($i\in J'$, $j\in J''$), donc $J$ ne serait pas un intervalle.

% archive sheet number: 269

\scanpage{271}
\mspage{12}

\uline{Corollaire} L'ensemble des \deleted{\unsure{points}}\unsure{sommets} critiques de $I$ \deleted{\unsure{déf}} \unsure{est|peut} aussi l'ens. des sommets du \deleted{\unsure{la}} plus petit sous-\deleted{catégorie pleine \uline{connexe}} \inserted{intervalle} de $I$ contenant les extrémités.

Plus généralement, si $i<j$ sont deux sommets de $I$, le plus petit \inserted{ss-} \deleted{catégorie pleine} intervalle de $I$ contenant $i,j$ est formé de $i,j$ et des sommets critiques $s$ tels que $i<s<j$.

\uline{Définition} Un intervalle est dit \uline{réduit} \inserted{(ou \uline{\unsure{irrédondant|irrédundant}})} si tous ses sommets sont critiques, ou ce qui revient au même, si ses chaînes structurales sont de longueur 1.

Un intervalle orienté \uline{réduit} est connu, à isom unique près, quand on connaît

a) Sa longueur $n$

b), Sa direction initiale $+$ (directe) ou $-$ (rétrograde)

\uline{NB} Cette direction est la $\mathord{=}$ que que la direction finale, \deleted{\unsure{Dans le cas d}} si $n$ est impaire, elle est opposée si $n$ est paire.

Ainsi, suivant les deux cas \inserted{de direction initiale}, on trouve

\diagram{\begin{tikzcd}[ampersand replacement=\&,column sep=small]
0\arrow[r] \& 1 \& 2\arrow[l]\arrow[r] \& \ldots \& \ldots\,\unsure{\bullet|\cdot}\, n \& \left(\begin{gathered}\text{Avec }n-1\longrightarrow n\quad\text{si }n\text{ imp.}\\\phantom{\text{Avec }}n-1\longleftarrow n\quad\text{si }n\text{ pair}\end{gathered}\right)
\end{tikzcd}}

\diagram{\begin{tikzcd}[ampersand replacement=\&,column sep=small]
0 \& 1\arrow[l]\arrow[r] \& 2 \& \ldots\arrow[l] \& \unsure{\to|\cdot} \& \ldots\,n \& \left(\begin{gathered}\text{Avec }n-1\longleftarrow n\quad\text{si }n\text{ imp}\\\phantom{\text{Avec }}n-1\longrightarrow n\quad\text{si }n\text{ pair}\end{gathered}\right)
\end{tikzcd}}

\deleted{\uline{\unsure{Définition}}} \inserted{Prop. 6} \deleted{Le \deleted{plus} sous-intervalle de $I$ formé des sommets critiques est \deleted{appelé l'intervalle réduit} \inserted{non seulement le plus petit} sous-intervalle contenant les extrémités, mais aussi le plus grand sous-intervalle irrédundant.}

\deleted{\unsure{Car :}} un sous-intervalle $J$, ayant $i$ et $j$ ($i<j$) ``comme'' plus petit et plus grand sommet, est \deleted{irrédundant} \inserted{réduit} sss \deleted{\unsure{les som}} tout \deleted{\unsure{$j\in$}} $s\in J$ distinct de $i$ et de $j$ est critique \unsure{\}|?}.

\uline{Cor.} La ss-cat. pleine de $I$ formée des sommets critiques est un sous-intervalle. On l'appelle l'intervalle réduit \uline{associé} à $I$, noté $I_{\mathrm{red}}$.

% archive sheet number: 270

\scanpage{272}
\mspage{13}

3. \deleted{\uline{\unsure{Composition} des intervalles orientés}} \uline{Chemins}. \uline{Catégorie des chemins d'un type donné}.

Soient $I$ un intervalle orienté, et $X$ une catégorie. On appelle \uline{chemin de type} $I$ dans $X$ \unsure{tout} foncteur $I\unsure{\longrightarrow|\text{---}}X$, \deleted{\unsure{et le}} et \uline{catégorie des chemins de type} $I$ \uline{dans} $X$ la catégorie

\[\underline{\mathrm{Ch}}_I(X)=\underline{\mathrm{Hom}}(I,X).\]

Comme $I$ est \uline{canoniquement} isom. à une catégorie intervalle standard $I_0$, \deleted{le ``type''} \deleted{\unsure{les données}} on peut identifier les chemins de type $I$, ou de type $I_0$. Le type \inserted{$T$} d'un chemin \inserted{de longueur $n$ ($=\operatorname{card}I$)} dans $X$ peut donc être décrit par une suite de $n$ signes $+$ ou $-$

\[\unsure{\tau|T|\xi|\varepsilon}_1,\tau_2,\ldots,\tau_n\qquad(\tau_i\in\{-1,+1\}\text{ pour }1\leq i\leq n).\]

À cause du fait que $I$ est engendrée par ses flèches de transition, un foncteur $c:I\unsure{-|\longrightarrow}X$ est dét\inserted{\unsure{er}}miné \deleted{de façon univ} quand on connaît

a) les $c(i)$ ($i\in I$) i.e une suite de \inserted{$n+1$} objets de $X$ :

\[\deleted{\unsure{a|x}}A_0,\deleted{a}A_1,\ldots,\deleted{a}A_n\,;\]

% slanted
\margin{$\unsure{A|a}_0$ est dit la \uline{source} \unsure{du chemin|chemin}, $a_n$ est son \uline{but}.}

b) Les \deleted{\unsure{fl}} $c(\underline{u}_i)$ (où les $\underline{u}_i$ sont les flèches de transition de $I$) i.e. pour tout $i\unsure{\notin|\in}[1,n]$ une flèche \inserted{$u_i$} dans $X$ entre $\deleted{a}A_{i-1}$ et $\deleted{a}A_i$, allant de $\deleted{a}A_{i-1}$ dans $\deleted{a}A_i$ si $\underline{u}_i$ est directe ($\tau_i=+1$), et de $\deleted{a}A_i$ dans $\deleted{a}A_{i-1}$ si $\underline{u}_i$ est rétrograde ($\tau_i=-1$),

et \deleted{aux} \inserted{à toute telle} donnée a), b) correspond un foncteur. Ainsi on pourra identifier les chemins de type $\deleted{\unsure{I}}\unsure{T|I}$

% archive sheet number: 271

\scanpage{273}
\mspage{14}

dans $X$ aux diagrammes du type précisé de $X$. \deleted{\unsure{E}}Et les \deleted{\unsure{homo}}morphismes entre chemins correspondent de $=$ aux morphismes entre diagrammes $A$ et $\deleted{\unsure{a}}A'$, une flèche $A\xrightarrow{\deleted{\unsure{f|g}}\alpha}\deleted{\unsure{a}}A$ étant donc \deleted{\unsure{une}} une suite de \inserted{$n+1$} flèches

\[
\alpha_i\colon A_i\mathrel{\unsure{-|\to}}A'_i,\quad i\unsure{\vdash|\in}[0,n],
\]

satisfaisant la condition de commutativité pour les $n$ carrés

\diagram{\begin{tikzcd}[ampersand replacement=\&]
A_{i-1}\arrow[r,"u_i"]\arrow[d,"\alpha_{i-1}"'] \& A_i\arrow[d,"\alpha_i"] \& {}\arrow[d,phantom,"{\text{(cas }\underline{u}_i\text{ directe)}}" description] \\
A'_{i-1}\arrow[r,"u'_i"] \& A'_i \& {}
\end{tikzcd}}

ou

\diagram{\begin{tikzcd}[ampersand replacement=\&]
A_{i-1}\arrow[d,"\alpha_{i-1}"'] \& A_i\arrow[l,"u_i"']\arrow[d,"\alpha_i"] \& {}\arrow[d,phantom,"{\text{(cas }\underline{u}_i\text{ rétrogr.)}}" description] \\
A'_{i-1} \& A'_i\arrow[l,"u'_i"'] \& {}
\end{tikzcd}}

Si on a une flèche entre deux sommets quelconques $i,j$ de $I$, avec $i<j$ et $j$ non \unsure{success.|succ.} immédiat de $i$, alors $i,j$ appartiennent à un $=$ intervalle critique, et dans $I$ on aura

\diagram{\begin{tikzcd}[ampersand replacement=\&]
i\arrow[r,"\underline{u}_{i+1}"] \& i+1\arrow[r,"\underline{u}_{i+1}"] \& \cdots \& j-1\arrow[r,"\underline{u}_j"] \& j
\end{tikzcd}}

une chaine de \deleted{\unsure{critiques}} longueur $j-i\geq 2$

ou

\diagram{\begin{tikzcd}[ampersand replacement=\&]
i \& i+1\arrow[l] \& \cdots\arrow[l] \& j-\unsure{1|j} \& j\arrow[l]
\end{tikzcd}}

et la flèche correspondante entre $A_i$ et $A_j$ s'obtient en composant les flèches $u_i$ correspondantes.

% archive sheet number: 272

\scanpage{274}
\mspage{15}

\deleted{Ceci rend immédiate la}

\deleted{\uline{Proposition 7} Soit $s\in I$ un élément qui}

\uline{Exemples}

1. $I=\Delta_0=e$, $\underline{\mathrm{Ch}}_I(X)\simeq X$. Les chemins de longueur $0$ s'identifient aux objets \inserted{$A$} de $X$. L'unique élément $\mathrm{id}_e$ de $\mathrm{Fl}(I)$ \deleted{\unsure{dev.|dev}} donne $\mathrm{id}_A$ dans $X$.

2. $I=\mathbb{\Delta}_1^+$, $\underline{\mathrm{Ch}}_I(X)\simeq\underline{\mathrm{Fl}}(X)$. Les chemins de \deleted{longueur $1$} type $\mathbb{\Delta}_1^+$ sont \inserted{ceux de longueur $1$ et \uline{directs}}, i.e. donnés par une flèche

\diagram{\begin{tikzcd}[ampersand replacement=\&]
a_0\arrow[r] \& a_1
\end{tikzcd}}

\deleted{qui est \unsure{la|le}} \inserted{dont la} source est aussi la source $a_0$ du chemin. Par contre, \deleted{\unsure{p.}} les chemins de type $\mathbb{\Delta}_1^-$ sont donnés par des diagrammes

\diagram{\begin{tikzcd}[ampersand replacement=\&]
a_0 \& a_1\arrow[l]
\end{tikzcd}}

i.e par des flèches dont le but est considéré comme source du chemin. Bien sûr,

% slanted
\margin{\unsure{?}}

\deleted{\unsure{Ch}} les catégories $\underline{\mathrm{Ch}}_{\mathbb{\Delta}_{\unsure{1}}^{+}}(X)$ et $\underline{\mathrm{Ch}}_{\mathbb{\Delta}_{\unsure{1}}^{-}}(X)$ sont canoniquement isomorphes, mais par un isomorphisme qui \uline{échange source et but \unsure{du|des} \unsure{chemins|chemin}}. Nous nous garderons donc bien de confondre les deux chemins précédents.

% archive sheet number: 273
